\section{Introduction}
\label{sec:introduction}

Black-box optimization is used when objective values can be obtained by evaluation but an analytic model, derivative information, or a tractable problem structure is unavailable. In this setting, the choice of optimizer matters: algorithms expose complementary behavior across problem instances, while no member of a practical portfolio is uniformly preferable. Automated algorithm selection addresses this heterogeneity by mapping an instance representation to a portfolio decision. A common workflow evaluates a space-filling design, computes exploratory landscape analysis (ELA) features, and then applies a supervised selector \cite{kerschkeAutomatedAlgorithmSelection2019,guoAutomatedAlgorithmSelection2025}. This workflow makes landscape analysis a prerequisite for selection.

That prerequisite is not innocuous. Landscape characterization comprises families of measures with different assumptions, sampling requirements, search dependence, and computational demands \cite{malanSurveyTechniquesCharacterising2013,pragerPflaccoFeaturebasedLandscape2024}. The objective evaluations used to construct a landscape sample consume budget that could otherwise be assigned to optimization, and the appropriate feature-computation fraction depends on the selection scenario \cite{blomInfluenceFeatureComputation2026}. Contemporary systems therefore count feature evaluations or omit expensive groups under a restricted budget \cite{guoAutomatedAlgorithmSelection2025}. An ELA-informed portfolio decision may be useful and still fail to justify the resources spent to obtain its descriptors at a particular search state.

At the same time, an optimization prefix already produces information. Search-trajectory networks characterize convergence patterns, recurrent transitions, and attraction regions from ordinary optimizer runs \cite{ochoaSearchTrajectoryNetworks2021}. DynamoRep preserves the longitudinal evolution of population summaries, while Opt2Vec learns from populations visited during search \cite{cenikjDynamoRepTrajectoryBasedPopulation2023,korosecOpt2VecContinuousOptimization2024}. Probing trajectories use short objective-value sequences for a subsequent portfolio choice, and later work shows that their evaluation depends on both the classifier and the problem-level split \cite{renauUtilityProbingTrajectories2024,renauAlgorithmSelectionProbing2025b}. Reusing evaluated points avoids a separate objective-function sample, but representation and inference still have computational cost. More importantly, problem classification or direct algorithm selection does not establish that a pre-query behavior summary predicts the net value of acquiring an independent descriptor query.

Generalization is a second obstacle. Algorithm-selection performance can appear strong when training and evaluation sets contain closely related instances. Under demanding problem and problem-combination splits, several classical and learned landscape representations did not improve over a single-best-solver baseline in the evaluated affine-BBOB setting \cite{cenikjLandscapeFeaturesSingleobjective2025}. The closest per-run warm-starting study likewise reported that results obtained on BBOB did not transfer directly to YABBOB \cite{kostovskaPerrunAlgorithmSelection2022}. Recent reviews identify heavy reliance on BBOB, heterogeneous feature protocols, and limited cross-benchmark evidence as continuing concerns \cite{cenikjSurveyFeaturesUsed2026}. Random instance partitions and within-training scores are therefore insufficient evidence for transfer to unseen function families or benchmark suites.

This paper studies a decision that precedes feature acquisition: given an optimization prefix and its algorithm-agnostic behavioral history, should a fixed landscape-descriptor query be executed at the current state? We call this setting \emph{Decision-before-Feature}. It differs from trajectory-based per-run selection, which constructs a trajectory representation and chooses a second-stage algorithm \cite{jankovicTrajectorybasedAlgorithmSelection2022,kostovskaPerrunAlgorithmSelection2022}, and from dynamic algorithm selection, which repeatedly changes search actions during a run \cite{guoDeepReinforcementLearning2024b}. It also differs from landscape-aware operator or population control \cite{sallamLandscapebasedAdaptiveOperator2017,zhouAdaptiveMultipopulationArtificial2024}. The object here is not a parameter, operator, or optimizer update rule, but whether to acquire additional information before a downstream portfolio decision.

We formulate query value as the difference in clipped terminal $\log_{10}$ gap between a native no-query continuation and a query-plus-selector path from the same complete optimizer state, less a weighted $\log_{10}$ ratio of their state-to-terminal future-path runtimes. The primary time for each path is the median of three prespecified censored replays: a completed replay retains its observed time, whereas a timed-out or failed replay contributes the larger of its observed time and role-specific timeout. The already executed shared prefix is a sunk cost; full FE-zero-to-terminal wall-clock remains a separate policy endpoint. Memory is measured and reported separately in the primary analysis, for which $\lambda_M=0$. The query evaluations already reduce the remaining optimization budget and are not deducted twice. Query-sample evaluations do not enter the optimizer population, but their best point and observed target hit enter the operational Query endpoint. We distinguish an observed target hit, path completion, and completion-conditioned endpoint success; standard ERT uses the observed hit even when a later path failure occurs. Offline labels use fold-specific downstream selectors, whereas the deployed gate observes only behavior available before the query. Query descriptors, function identity, algorithm identity, optimizer-specific parameters, and known-optimum gaps are excluded from that gate. The resulting question is not whether landscape features can sometimes support selection, but when the expected downstream difference of the evaluated query is sufficient to justify its acquisition cost.

The study is organized around five research questions. RQ1 characterizes the statewise utility of the primary \query{} query. RQ2 tests whether algorithm-agnostic behavior adds decision value beyond elapsed budget and a dimension-stratified time-only comparator. RQ3 compares the resulting policy with fixed, random, traditional landscape-based, behavior-only, and portfolio references under common resource accounting. RQ4 reports the inspected six-function BBOB-validation set and CEC2017 as development evaluations, and reserves CEC2022 and engineering problems for prospective evaluation after their complete configurations and endpoints have been fixed. Each suite is reported separately. The BBOB-validation estimand is a finite-set mean, not a function-superpopulation or independent-confirmation estimand. RQ5 examines the six prespecified T0/B1/B2/B2+Motion/B2+Maturity/B3 groups, model-appropriate coefficient or discriminant stability, maturity--utility associations, and variation across function, dimension, and query configuration. RQ4 and RQ5 remain empirical questions; their intended evaluations are not evidence that transfer or explanatory stability has already been obtained.

The contributions are methodological. First, we define execution of a fixed landscape-descriptor query as a resource-aware prediction problem with a state-conditional net-utility target. Second, we construct the gate input from permutation-invariant native-update histories while excluding algorithm identity and query descriptors. Third, we separate this gate from a query-specific downstream action-loss selector and make native continuation and population-transfer initialization explicit. Fourth, we specify nested grouped-by-function out-of-fold model selection and train-only threshold fitting, together with an evaluation structure connecting decision utility, final optimization loss, query use, and measured resource cost. The grouping unit is function ID, not a claim of generalization across a classical landscape-family taxonomy. Among the literature reviewed here, we did not identify the same combination of pre-query information timing, matched-state utility labels, and an independently acquired query. Whether the method yields a predictive or resource advantage is determined only by RQ1--RQ5.

The remainder of the paper reviews landscape-based and trajectory-based selection, behavior characterization, adaptive optimization, and generalization; formalizes the decision and utility target; describes the offline construction and learning procedure; and presents the RQ-aligned evaluation and reproducibility design.
